
Donc il suffit de montrer que l'on a pour tout \(k>1\)
\[
 \det(\gamma^k(E-\underline{n}))=1,
\]
où, comme ci-dessus, \(E=\cl^{\bullet}(\underline E)\), \(\underline{n}=\epsilon(E)\). Or, on a
\[
 \gamma^k(x)=\lambda^k(x+k-1)
\]
pour tout \(x\in K^{\bullet}(X)\), de sorte que
\[
 \det(\gamma^k(E-\underline{n}))
 =
 \det(\lambda^k(E-\underline{n}+k-1))
 =
 \det\left(\sum_{p+q=k}\binom{-\underline{n}+k-1}{p}\lambda^q(E)\right).
\]
Mais (6.2.10) et (6.5.3) donnent
\[
 \det(\lambda^q(E))=(\det(E))^{\binom{\underline{n}-1}{\underline{q}-1}}
\]
pour tout \(q>0\), d'où
\[
 \det(\gamma^k(E-\underline{n}))=(\det(E))^{\underline{r}},
\]
avec
\[
 \underline{r}=\sum_{\substack{p+q=k\\ q\geq1}}
 \binom{-\underline{n}+k-1}{p}\binom{\underline{n}-1}{q-1}.
\]
Or, \(\underline{r}\) est une section de \(\mathbb Z\), partout égale au coefficient de \(t^{k-1}\) dans \((1+t)^{k-2}\), donc nulle. On a gagné.

\section*{6. Appendice : Calcul des déterminants des faisceaux localement libres}

\subsection*{6.1.}
Considérons d'abord un topos annelé quelconque \(X=({C},\underline{O})\). Pour toute section positive \(\underline{n}\in\Gamma(\mathbb Z_X)\) du faisceau constant \(\mathbb Z\) et pour tout \(\underline{O}\)-module \(M\in\Mod_{\underline{O}}(X)\), il est clair comment définir les \(\underline{O}\)-modules
